\chapter{The soft PLC}
\label{ch:plc}

A programmable logic controller is a computer built to run one program
forever, deterministically, on a fixed cycle. This chapter describes the PLC
model pyplant implements: the scan cycle, the program and its precedence
rules, protective logic, timers, the drift-free scheduler and the
communication layers.

\section{The scan cycle}

Every PLC scan does three things, always in this order:

\begin{enumerate}
  \item \textbf{Read inputs.} All field inputs are copied into an
        \emph{input image table} in one operation.
  \item \textbf{Execute logic.} The program runs against the frozen image.
        It never reads a field value directly, so every rung in one scan sees
        the same consistent snapshot.
  \item \textbf{Write outputs.} The \emph{output image table} is written to
        the field in one operation.
\end{enumerate}

\begin{figure}[htbp]
\centering
\scalebox{0.8}{%
\begin{tikzpicture}[node distance=8mm and 10mm]
  \node[term] (tick) {Ticker fires (every $T_s$)};
  \node[box,below=of tick,minimum width=52mm] (read) {Read inputs over \code{IOLink}\\$\to$ \code{PlcInput} image};
  \node[hot,below=of read,minimum width=52mm] (exec) {\code{Program.execute(image)}\\alarms, interlock, PID};
  \node[box,below=of exec,minimum width=52mm] (write) {Write \code{PlcOutput} image\\over \code{IOLink}};
  \node[box,below=of write,minimum width=52mm] (diag) {Publish diagnostics\\\tg{SCAN_MS}, \tg{SCAN_OVERRUNS}, \tg{PLC_HEARTBEAT}};
  \node[decision,right=18mm of exec] (late) {late by\\$\ge T_s$?};
  \node[bad,right=of late] (ovr) {count\\overrun};
  \draw[flow] (tick) -- (read);
  \draw[flow] (read) -- (exec);
  \draw[flow] (exec) -- (write);
  \draw[flow] (write) -- (diag);
  \draw[flow] (diag.west) -- ++(-10mm,0) |- (tick.west);
  \draw[dflow] (diag.east) -| (late.south);
  \draw[flow] (late) -- node[lbl,above]{yes} (ovr);
\end{tikzpicture}}
\caption{One scan of the PLC engine (\code{pyplant/plc/scan.py}).}
\label{fig:scan}
\end{figure}

The engine (\code{ScanEngine}) does not know whether its I/O is in-process or
Modbus: all I/O goes through the \code{IOLink} protocol. The program
(\code{Program}) does not know about I/O at all: \code{execute()} is a pure
function of the input image and the program's retained state, which is what
makes it testable scan by scan.

\section{The program}

\subsection{Precedence}
Highest first, as on a real unit:
\begin{enumerate}
  \item \textbf{Interlock.} Latched on a trip condition. Forces the pump off,
        drives the valve shut, and keeps the PID tracking 0\% so a reset
        cannot slam the valve open. Clears only when the operator resets
        \emph{and} the condition has gone.
  \item \textbf{Manual.} The operator's output demand passes through; the PID
        tracks it so that returning to auto is bumpless.
  \item \textbf{Auto.} The PID holds \tg{LT101} at \tg{LIC101_SP}.
\end{enumerate}

\begin{figure}[htbp]
\centering
\scalebox{0.8}{%
\begin{tikzpicture}[node distance=7mm and 12mm]
  \node[term] (in) {input image};
  \node[box,below=of in] (alm) {evaluate alarms\\with hysteresis};
  \node[decision,below=of alm] (trip) {level $>$ 550 mm or\\low flow for 15 s?};
  \node[bad,right=of trip] (latch) {latch\\interlock};
  \node[decision,below=of trip] (rst) {latched and\\reset pressed?};
  \node[good,right=of rst] (unl) {unlatch,\\reset timer};
  \node[decision,below=of rst] (il) {interlocked?};
  \node[bad,left=of il] (safe) {OP = 0\%, pump off,\\PID tracks 0\%};
  \node[decision,below=of il] (auto) {auto?};
  \node[hot,below left=7mm and 2mm of auto] (pid) {PID step\\eq.~\eqref{eq:velocity}};
  \node[box,below right=7mm and 2mm of auto] (man) {OP = manual demand,\\PID tracks it};
  \node[term,below=22mm of auto] (out) {output image};
  \draw[flow] (in) -- (alm);
  \draw[flow] (alm) -- (trip);
  \draw[flow] (trip) -- node[lbl,above]{yes} (latch);
  \draw[flow] (trip) -- node[lbl,right]{no} (rst);
  \draw[flow] (rst) -- node[lbl,above]{yes} (unl);
  \draw[flow] (rst) -- node[lbl,right]{no} (il);
  \draw[flow] (unl) |- (il);
  \draw[flow] (latch.east) -- ++(6mm,0) |- (il.east);
  \draw[flow] (il) -- node[lbl,above]{yes} (safe);
  \draw[flow] (il) -- node[lbl,right]{no} (auto);
  \draw[flow] (auto) -| node[lbl,pos=0.3,above]{yes} (pid);
  \draw[flow] (auto) -| node[lbl,pos=0.3,above]{no} (man);
  \draw[flow] (pid) |- (out);
  \draw[flow] (man) |- (out);
  \draw[flow] (safe.west) -- ++(-5mm,0) |- (out.west);
\end{tikzpicture}}
\caption{Decision flow of \code{Program.execute()} in one scan.}
\label{fig:program}
\end{figure}

\subsection{Protective functions}

\begin{center}
\begin{tabularx}{\textwidth}{@{}lXX@{}}
\toprule
Function & Trigger & Action \\
\midrule
\tg{ALM_LT101_HI} & \tg{LT101} $>$ 500\,mm & alarm, priority medium \\
\tg{ALM_LT101_LOLO} & \tg{LT101} $<$ 50\,mm & alarm, priority high \\
\tg{ALM_PT101_HIHI} & \tg{PT101} $>$ 58\,kPa & alarm, priority high \\
\tg{ILK_PMP101} & \tg{LT101} $>$ 550\,mm, or \tg{FT101} $<$ 1\,L/min for 15\,s while running & latched trip: pump off, valve shut; priority critical \\
\bottomrule
\end{tabularx}
\end{center}

All thresholds are read from the tag dictionary, so a respanned instrument
cannot leave a stale limit behind in the logic.

\begin{theory}[Why minimum flow and not high pressure]
A centrifugal pump against a shut valve rises to its shut-off head, 60\,kPa
here. A high-pressure trip below that would fire during any legitimate
low-output excursion. The real hazard of dead-heading is that no flow carries
heat away, so the trip is \emph{low flow on a time delay} and high pressure
is only an alarm.
\end{theory}

\subsection{Hysteresis}
A comparator that switches exactly at a limit chatters when the signal rests
on it, and an alarm list that chatters is one operators stop reading. Each
alarm is a latching comparator: a high alarm sets when $x \ge L$ and clears
only when $x \le L - \delta$, with deadbands
$\delta = 10$\,mm for levels, 2\,kPa for pressure and 0.5\,L/min for flow.

\begin{figure}[htbp]
\centering
\begin{tikzpicture}
\begin{axis}[width=0.58\textwidth,height=3.5cm,xlabel={time},ylabel={level (mm)},
  xmin=0,xmax=10,ymin=470,ymax=520,ytick={490,500},yticklabels={$L-\delta$,$L$},
  xtick=\empty,grid=major,grid style={linegrey}]
\addplot[accent,thick,smooth] coordinates {(0,480)(2,497)(3,502)(4,498)(5,503)(6,496)(7,492)(8,488)(10,485)};
\addplot[badred,very thick,const plot] coordinates {(0,474)(2.6,474)(2.6,514)(7.5,514)(7.5,474)(10,474)};
\node[font=\footnotesize,badred] at (axis cs:5,517) {alarm on};
\end{axis}
\end{tikzpicture}
\caption{Hysteresis: the alarm (red) sets at $L$ and clears only below
$L-\delta$, so the wobble between 496 and 503\,mm produces a single alarm.}
\label{fig:hyst}
\end{figure}

\subsection{Timers count scans}
The minimum-flow trip uses an on-delay timer (TON): its output goes true once
its input has been true for the preset, and resets the moment the input goes
false. The timer counts \emph{scans}, not seconds. Accumulating
\code{elapsed += 0.1} reaches $0.9999999999999999$ after ten scans, and a
timer compared with \code{>=} against a 1.0\,s preset would never fire. Real
PLC timers count ticks for exactly this reason. The preset is converted once
to $\operatorname{round}(T/T_s)$ scans.

\section{Drift-free scheduling}

\code{await asyncio.sleep(period)} in a loop schedules the next wake relative
to when the previous one \emph{finished}, so every microsecond of work is
added permanently to the phase. The \code{Ticker} in \code{pyplant/clock.py}
computes absolute targets from a fixed origin instead:
\begin{equation}
  t_n = t_0 + n\,\frac{T}{\text{speed}} .
\end{equation}
Work eats into the slack, never into the phase. With a 70\,ms body on a
100\,ms period, 100 ticks land at exactly 10.000\,s rather than 17.000\,s; a
test asserts this to $10^{-12}$\,s.

\begin{figure}[htbp]
\centering
\scalebox{0.8}{%
\begin{tikzpicture}[x=1.1cm,y=1cm]
  \draw[->] (0,0) -- (10.5,0) node[right]{time};
  \foreach \i in {0,...,9}{\draw (\i,0.08) -- (\i,-0.08); }
  \node[left] at (0,0.6) {\small drift-free};
  \foreach \i in {0,...,8}{\fill[accent!60] (\i,0.35) rectangle (\i+0.7,0.85);}
  \node[left] at (0,-0.8) {\small sleep(period)};
  \foreach \i in {0,...,5}{\fill[badred!50] (\i*1.7,-0.55) rectangle (\i*1.7+0.7,-1.05);}
  \node[font=\footnotesize] at (5,1.25) {each tick starts on the grid; work uses slack};
  \node[font=\footnotesize] at (5,-1.45) {each tick starts after the last one ends: phase slips by the work time};
\end{tikzpicture}}
\caption{Absolute scheduling (top) against relative sleeping (bottom) with a
70\% duty cycle.}
\label{fig:ticker}
\end{figure}

\subsection{The engine measures itself}
\tg{SCAN_MS}, \tg{SCAN_OVERRUNS} and \tg{PLC_HEARTBEAT} are real tags, trended
like any process value. The published statistics are always the previous
scan's, because this scan's duration is unknown until its write completes,
which is what real controllers do. When the engine cannot keep up, lateness
grows without bound against the fixed grid while the overrun count increments
once per scan; lateness is the signal to watch.

\section{The tag dictionary}

\code{pyplant/tags.py} is the single source of truth for every signal. Each
\code{Tag} carries a description, kind (analog or bool), owning process,
Modbus table and address, engineering unit and span, integer scale factor,
alarm limits, writability and the OPC UA area. Every other module imports its
I/O definitions from it; nothing else may hard-code an address or a scale
factor. Appendix~\ref{app:tags} lists all tags.

\subsection{Wire representation}
Analog values travel in 16-bit Modbus holding registers as
\begin{equation}
  w = \operatorname{round}(v\cdot s), \qquad
  w \in [-32768, 32767],
\end{equation}
sent as an unsigned two's-complement word. The scale $s$ is chosen per tag so
the span fits; this is checked at import. Out-of-range values are clamped,
not wrapped, so a bad value reads as a pegged instrument rather than a sign
flip. For \tg{LT101}, $s = 10$: 432.7\,mm travels as 4327 with a resolution
of 0.1\,mm. Booleans travel in coils as 0 or 1.

\subsection{Ownership and blocks}
Each tag's \code{owner} (\code{PLANT}, \code{PLC}, \code{OPER}, \code{ML})
is the only process that may write it. Addresses are grouped into declared
\emph{blocks} so a client polls a whole span in one Modbus transaction, and
import fails if a block has a gap, two blocks overlap, or a tag sits in a
block owned by another process.

\begin{theory}[No block may be written by two parties]
If one block held both the operator's setpoint and the PLC's readback, an
operator write arriving between the PLC's read and its publish would be
overwritten 100\,ms later. The faceplate is therefore split into
\code{OPER_SP}/\code{OPER_CMD} (HMI writes, PLC reads) and
\code{PLC_FACE}/\code{PLC_DIAG}/\code{PLC_ALARMS} (PLC writes, HMI reads).
\end{theory}

\section{Modbus TCP: the field layer}

\begin{figure}[htbp]
\centering
\scalebox{0.8}{%
\begin{tikzpicture}[node distance=6mm and 22mm]
  \node[box,minimum width=36mm,minimum height=30mm] (plant) {\textbf{plant node}\\Modbus slave\\port 5400\\[3pt]\code{PLANT_MEAS} HR 0--5\\\code{PLANT_STATUS} coil 0\\\code{PLC_CMD} HR 10\\\code{PLC_CMD_BITS} coil 10};
  \node[hot,minimum width=36mm,minimum height=30mm,right=of plant] (plc) {\textbf{PLC node}\\client below,\\server above\\port 5402};
  \node[box,minimum width=30mm,right=of plc] (scada) {SCADA clients\\(Modbus or\\OPC UA 5403)};
  \draw[flow] ([yshift=8mm]plc.west) -- node[lbl,above]{1. read meas} ([yshift=8mm]plant.east);
  \draw[flow] ([yshift=2.5mm]plc.west) -- node[lbl,above]{2. read status} ([yshift=2.5mm]plant.east);
  \draw[flow] ([yshift=-3mm]plc.west) -- node[lbl,above]{3. write OP} ([yshift=-3mm]plant.east);
  \draw[flow] ([yshift=-8.5mm]plc.west) -- node[lbl,above]{4. write pump} ([yshift=-8.5mm]plant.east);
  \draw[flow,<->] (plc) -- (scada);
\end{tikzpicture}}
\caption{Four Modbus round trips per scan, measured at 0.7\,ms against a
100\,ms budget.}
\label{fig:modbus}
\end{figure}

The plant is a Modbus \emph{slave}, as a field device is. It has no controller,
no alarms and no memory of intent: it applies whatever is in its command
registers, so if the PLC stops writing, the valve holds its last position, as
a real one does. The PLC is both client (towards the plant) and server
(towards SCADA). It mirrors the plant measurements into its own store so the
HMI polls one address instead of two. Reads are by block, never by tag: one
\code{read_holding_registers} for six measurements is one round trip.

\section{OPC UA: the information layer}

OPC UA is the boundary between registers and meaning. Below it everything is
addresses and scale factors; above it the historian and HMI speak tag names
with units attached and never import a Modbus symbol.

The address space is generated from the tag dictionary. Each tag's
\code{area} becomes its folder path, e.g.\ \code{Plant/Unit1/LT101},
\code{PLC/LIC101/LIC101_SP} (writable) or \code{PLC/Diagnostics/SCAN_MS}.
Every analog node carries \code{EngineeringUnits} and \code{EURange} from its
declared span, so a client can scale a trend without being told anything about
the plant.

The server is a \emph{view} of the same \code{TagStore} the Modbus server
holds, not a second process image. Values are pushed outward on a refresh tick
only for nodes that changed; client writes are pushed back into the store. To
stop a refresh from reverting a fresh operator write, \code{refresh()} tracks
the last value each writable tag was seen to hold and overwrites a writable
node only when the store changed underneath it.

Clients use \emph{subscriptions}: the server pushes changes at a publishing
interval, which is why the historian is not a poller.
